\documentclass[11pt,letterpaper]{article}
\usepackage[T1]{fontenc}
\usepackage[utf8]{inputenc}
\usepackage{lmodern,microtype}
\usepackage[margin=1in]{geometry}
\usepackage{amsmath,amssymb,amsthm,mathtools}
\usepackage{booktabs,array,longtable,enumitem}
\usepackage[dvipsnames]{xcolor}
\usepackage{listings}
\usepackage{hyperref}
\hypersetup{colorlinks=true,linkcolor=MidnightBlue,urlcolor=MidnightBlue,citecolor=MidnightBlue,
 pdftitle={The fundamental theorem of algebra in intuitionistic primitive recursive arithmetic},
 pdfauthor={Prepared for Vladimir Reshetnikov}}
\usepackage{fancyhdr}
\pagestyle{fancy}\fancyhf{}
\fancyhead[L]{\small Constructive FTA in iPRA}
\fancyhead[R]{\small Explicit rational construction}
\fancyfoot[C]{\thepage}
\setlength{\headheight}{14pt}
\setlength{\emergencystretch}{3em}
\setlist{itemsep=3pt,topsep=5pt}
\lstset{basicstyle=\small\ttfamily,breaklines=true,columns=fullflexible,
 frame=single,rulecolor=\color{black!20},showstringspaces=false}
\newtheorem{theorem}{Theorem}[section]
\newtheorem{lemma}[theorem]{Lemma}
\newtheorem{proposition}[theorem]{Proposition}
\newtheorem{corollary}[theorem]{Corollary}
\theoremstyle{definition}\newtheorem{definition}[theorem]{Definition}
\theoremstyle{remark}\newtheorem{remark}[theorem]{Remark}
\newcommand{\ipra}{\mathsf{iPRA}}
\newcommand{\N}{\mathbb N}\newcommand{\Q}{\mathbb Q}
\newcommand{\GQ}{\mathbb Q(i)}\newcommand{\C}{\mathbb C}
\newcommand{\nn}[1]{\mathcal N(#1)}
\newcommand{\abs}[1]{\left|#1\right|}
\newcommand{\ceil}[1]{\left\lceil#1\right\rceil}
\newcommand{\bl}{\operatorname{bl}}
\newcommand{\PR}{\mathrm{PR}}
\newcommand{\re}{\operatorname{Re}}\newcommand{\im}{\operatorname{Im}}
\newcommand{\Coeff}{\mathsf{Coeff}}
\newcommand{\Name}{\mathsf{Name}}
\newcommand{\Apart}{\mathsf{Apart}}
\newcommand{\Root}{\mathsf{Root}}

\title{\textbf{The fundamental theorem of algebra\\in intuitionistic primitive recursive arithmetic}\\[5pt]
\large A complete human-readable proof by bounded rational descent\\and a blueprint for a deep Lean formalization}
\author{Prepared for Vladimir Reshetnikov}
\date{November 14, 2026}

\begin{document}
\maketitle
\begin{abstract}
We give an explicit primitive recursive construction of a complex root of a
positively nonconstant polynomial whose coefficients are specified by primitive
recursive fast Cauchy names. The upper bound on the degree is arbitrary. A
positive lower-bound certificate for any nonconstant coefficient suffices;
the highest listed coefficient need not be known to be nonzero.

The proof works with rational polynomials at every iteration. A finite
comparison of $2n+1$ scales isolates a Taylor term, and a bounded rational
direction calculation produces a uniform decrease in the residual. A separate
finite product estimate bounds all iterates. Sampling them at an explicit
primitive recursive schedule gives a coherent fast Cauchy root name, rather
than merely unrelated approximate zeros. Every induction can be applied to
a quantifier-free predicate on finite numerical data.

The theorem inside first-order $\ipra$ is stated precisely as a schema for
fixed primitive recursive coefficient programs, with a uniform finite-code
special case for rational polynomials. We distinguish this from a second-order
assertion about all complex numbers. The accompanying Lean file is an
\emph{uncompiled specification and formalization interface}, not a completed
formal proof. Its intended endpoint is a derivation in ProveIt's deeply
embedded \texttt{PRAAx} theory. Exact rational reference programs and test
results accompany the proof.
\end{abstract}

\newpage
\tableofcontents
\newpage

\section{What the assertion means inside \texorpdfstring{$\ipra$}{iPRA}}
\label{sec:statement}

The usual fundamental theorem of algebra says that a nonconstant complex
polynomial has a complex root. Two points must be made explicit before asking
for a proof in intuitionistic primitive recursive arithmetic.

First, $\ipra$ is first-order arithmetic. Its variables range over natural
numbers, not over arbitrary functions or complex numbers. We therefore work
with numerical codes for rational approximations, supplied by fixed primitive
recursive functions. Second, constructive nonconstancy requires positive
information. We supply a rational certificate that a positive-degree
coefficient is apart from zero. We do not decide equality of complex numbers.

\subsection{The underlying arithmetic and its representations}

Use intuitionistic first-order logic with equality, the defining equations
for all primitive recursive functions, successor separation and injectivity,
and induction for quantifier-free formulas. This is the presentation of
$\ipra$ in the inspected \texttt{ConciseIPRA/PR.lean} source~\cite{proveit-pr}.
Integers, rationals with positive denominator, Gaussian rationals, finite
tuples, and finite sequences have their usual natural-number encodings.
Their arithmetic and comparisons are primitive recursive. A bounded
quantifier is an abbreviation for a primitive recursive Boolean fold, so
induction on such a bounded numerical assertion is available as
quantifier-free induction in the full primitive recursive language.

For $z=x+iy\in\GQ$, put
\[
 \nn z=x^2+y^2.
\]
All executed tests use this rational squared norm. In the explanatory proof,
$|z|$ denotes its nonnegative square root. For example,
$|z|\le r$ with rational $r\ge0$ means exactly $\nn z\le r^2$.
No comparison of arbitrary real numbers is used by the algorithm.

\begin{definition}[Coefficient names and a finite apartness certificate]
Fix a primitive recursive function $A(i,s,\mathbf p)$ with Gaussian-rational
values. The natural tuple $\mathbf p$ denotes ordinary numerical parameters.
Write $a_i(s)=A(i,s,\mathbf p)$. For an integer $n\ge1$, the name hypothesis is
\begin{equation}
 |a_i(u)-a_i(v)|\le 2^{-u}+2^{-v}
 \quad(0\le i\le n;\ u,v\in\N).
 \tag{N}\label{eq:names}
\end{equation}
The input includes $k,h\in\N$ with $1\le k\le n$ and
\begin{equation}
 |a_k(h+2)|\ge 2^{-h}.
 \tag{A}\label{eq:apart}
\end{equation}
Condition \eqref{eq:apart} is a single rational comparison, namely
$\nn{a_k(h+2)}\ge 2^{-2h}$.
\end{definition}

The hypotheses imply, for every $s\ge h+2$,
\begin{equation}
 |a_k(s)|\ge 2^{-h}-2^{-h-2}-2^{-s}
                 \ge 2^{-h-1}=:\lambda.
 \label{eq:apart-tail}
\end{equation}
Thus all sufficiently precise sampled polynomials are nonconstant.
The index $n$ is only an upper bound on the degree. In particular, the
argument allows a polynomial with a certified linear coefficient and an
undecided quadratic coefficient.

\begin{theorem}[Arithmetized constructive FTA]
\label{thm:main}
For every fixed primitive recursive coefficient program $A$, there are
explicit primitive recursive functions
\[
 \alpha(n,k,h,\mathbf p,t)\in\GQ,
 \qquad M(n,k,h,\mathbf p,t)\in\N
\]
whose defining construction is given below. In $\ipra$, from
\eqref{eq:names}, \eqref{eq:apart}, and $1\le k\le n$, one can prove
\begin{align}
 |\alpha(u)-\alpha(t)|&\le 2^{-t-2} &&(u\ge t),
 \label{eq:main-cauchy}\\
 \left|\sum_{i=0}^n a_i(s)\alpha(t)^i\right|&\le 2^{-t}
                   &&(s\ge M(t)).
 \label{eq:main-root}
\end{align}
The suppressed arguments of $\alpha$ and $M$ are unchanged. Consequently
$\alpha$ is an explicit fast Cauchy name for a root of the polynomial with
the given coefficient names.
\end{theorem}

This is a genuine $\ipra$ theorem schema: after a fixed primitive recursive
program $A$ is substituted, every expression in the assertion is a
first-order arithmetic expression. The construction is uniform in
$n,k,h,\mathbf p,t$. It does \emph{not} use a total primitive recursive
universal evaluator for arbitrary primitive recursive program codes; such an
evaluator is not available. The proof also gives an oracle-relative finite
algorithm, but this does not add function quantification to $\ipra$.

For a finite code of a rational polynomial, use constant names. Its
nonconstant coefficient and a suitable $h$ can be found by bounded numerical
operations. This gives one uniform first-order theorem for all nonconstant
rational polynomials. For complex coefficients the displayed schema is the
precise internal replacement for unrestricted quantification over $\C$.

\subsection{Why mere negated constancy is not silently assumed}

An explicit positive lower bound for a coefficient produces
\eqref{eq:apart} after a known precision shift. Mere
$\neg(a_k=0)$ need not produce such a bound intuitionistically. The issue is
already present in $az-1$: a root gives a bound on $|z|$ and therefore a
positive lower bound on $|a|$. An assertion producing this root from only
$\neg(a=0)$ would supply a nonzero-to-apartness principle; applied to the
usual real coding of a binary sequence, it yields a Markov-type principle.
No such principle is among our axioms.

The positive formulation above is stronger than a monic-only theorem: any
positive-degree coefficient can carry the certificate, with no prior
determination of the exact degree.

\section{Finite arithmetic tools and radical notation}
\label{sec:finite}

We record the elementary tools used throughout. They require only finite
calculations and quantifier-free induction on coded finite data.

\begin{enumerate}
\item Gaussian-rational arithmetic gives
\[
 \nn{zw}=\nn z\,\nn w,
 \qquad
 \nn z\,\nn w-(\re z\re w+\im z\im w)^2
       =(\re z\im w-\im z\re w)^2.
\]
The second identity gives Cauchy--Schwarz and the triangle inequality.
In particular,
\[
 \nn{\sum_{r<m}v_r}\le m\sum_{r<m}\nn{v_r}.
\]
Bounds with rational nonnegative right-hand sides can be checked by squaring.

\item Finite powers, sums, products, binomial coefficients, polynomial
translation, and polynomial evaluation are primitive recursive. Their
identities are proved by induction on finite indices. For example,
\begin{equation}
 x^i-y^i=(x-y)\sum_{r<i}x^{i-1-r}y^r
 \label{eq:power-diff}
\end{equation}
supplies every polynomial continuity estimate needed later.

\item For $x\ge-1$, finite Bernoulli induction gives
$(1+x)^m\ge1+mx$ in the instances used below. Finite geometric sums are
evaluated by multiplication and cancellation; no infinite series theorem
is needed.

\item If $r\in[0,1]\cap\Q$, $t$ ordinary rational bisections on $[0,1]$
give a lower endpoint $\operatorname{SQ}_t(r)$ within $2^{-t}$ of $\sqrt r$.
For $r\in[1,2]$, use $[1,2]$. The invariant is
$\ell^2\le r\le u^2$, $\ell\ge0$, and $u-\ell=2^{-t}$.
The decision is the rational comparison of the midpoint's square with $r$.
These bisections define the square-root names used as explanatory notation.
Also
\begin{equation}
 |\sqrt r-\sqrt s|\le\sqrt{|r-s|}\qquad(r,s\ge0)
 \label{eq:sqrt-holder}
\end{equation}
follows by squaring and using the nonnegative cross term.
\end{enumerate}

Nested radicals in the next section abbreviate explicitly approximable
algebraic quantities. They do not assert a completeness axiom. Bisection
with the displayed precision bounds replaces them by rational intervals.
Identities about these quantities mean identities of their explicit names.
At any finite target accuracy they follow from finite squaring identities
and interval bounds.

For completeness, precision propagation through any finite radical
expression used here has an explicit rational algorithm. Suppose rational
bounds $|x|,|y|\le B$ and a reciprocal lower bound $|x|\ge m>0$ are
available. To obtain output error at most $\varepsilon\le1$, it suffices
to demand the following input errors (and, for a square root, the stated
additional local bisection error):
\begin{center}
\begin{tabular}{@{}ll@{}}\toprule
Operation & Sufficient error budget \\
\midrule
$x\pm y$ & Each input error $\le\varepsilon/2$ \\
$xy$ & Each input error $\le\min(1,\varepsilon/(2B+1))$ \\
$1/x$ & Input error $\le\min(m/2,\varepsilon m^2/2)$ \\
$\sqrt x$ & Input error $\le\varepsilon^2/4$; local error
$\le\varepsilon/2$ \\
\bottomrule\end{tabular}
\end{center}
For a square-root input, clamp its rational approximation to the known
nonnegative interval; this cannot increase its error. The multiplication
bound follows from $|xy-x'y'|\le(2B+1)\delta$ for input errors at most
$\delta\le1$, and the reciprocal bound from
$|x^{-1}-(x')^{-1}|\le2\delta/m^2$. The square-root rule follows from
\eqref{eq:sqrt-holder}. Propagate these budgets backwards through the
finite expression tree, then evaluate forwards. The required bisection
count is bounded by a binary digit count of the rational interval width
divided by the target error. Both passes are finite primitive recursive
calculations. Magnitude bounds and reciprocal lower bounds are explicit
in every use below: for example, the normalized denominator is at least
$1$, and a nonzero rational norm is bounded below by the reciprocal of a
power of two computed from its component denominators. Thus
``approximate the radicals'' never conceals an unbounded precision search.

There is a useful last step when an intermediate argument uses arbitrarily
accurate radical bounds but its conclusion is rational. If $r\in\Q$ and
$r\le 2^{-m}$ has been proved for every $m$ by such finite bounds, then
$r\le0$ follows constructively. Decide the rational comparison $r\le0$.
In the other case, writing $r=a/b$ with positive integers gives $r\ge1/b$;
the explicit choice $m=\bl(b)+1$ contradicts $r\le2^{-m}$.
This is bounded arithmetic and rational decidability, not Markov's
principle. It eliminates the radical shorthand at the rational endpoints
of the lemmas below.

\section{A bounded rational direction algorithm}
\label{sec:direction}

The descent step needs an approximate $k$th root of a unit direction. Using
FTA for this would be circular. We instead give a finite algebraic
construction from square-root bisection and a truncated binomial polynomial.

\begin{lemma}[Direction]
\label{lem:direction}
Given $c\in\GQ\setminus\{0\}$, $k\ge1$, and rational
$0<\varepsilon\le1/16$, a bounded primitive recursive procedure constructs
$w\in\GQ$ such that
\[
 |w|\le1,\qquad
 \left|w^k-\frac{c}{|c|}\right|\le\varepsilon.
\]
\end{lemma}

\subsection{Three square roots move the target near \texorpdfstring{$1$}{1}}

Write $c=x+iy$, let $m=\max(|x|,|y|)>0$, and set
\[
 X=x/m,\quad Y=y/m,\quad
 a=\frac{X}{\sqrt{X^2+Y^2}}.
\]
Here $1\le X^2+Y^2\le2$ and $-1\le a\le1$.
Let $\sigma=-1$ if $y<0$ and $\sigma=1$ otherwise, a rational decision.
Define
\begin{align*}
 s_1&=\sqrt{(1+a)/2},&
 s_2&=\sqrt{(1+s_1)/2},\\
 s_3&=\sqrt{(1+s_2)/2},&
 t_3&=\sqrt{(1-s_2)/2},\qquad v=s_3+i\sigma t_3.
\end{align*}
Direct squaring, with the nonnegative square roots, shows
\[
 |v|=1,\qquad v^8=c/|c|.
\]
For detail, put
$v_1=s_1+i\sigma\sqrt{(1-a)/2}$ and
$v_2=s_2+i\sigma\sqrt{(1-s_1)/2}$.
Then $v^2=v_2$, $v_2^2=v_1$, and $v_1^2=c/|c|$.
The sign of $y$ verifies the imaginary part in the last equality. For
$y=0$, the convention selects a half-angle with nonnegative real part,
taking the positive imaginary sign on the negative real axis.

Since $s_1\ge0$, $s_2\ge2/3$, and $s_3\ge9/10$ (square these rational
lower bounds),
\begin{equation}
 |v-1|^2=2-2s_3\le1/5<1/4.
 \label{eq:v-near-one}
\end{equation}

\subsection{Computing a rational approximation}

Choose $\ell\ge3$, put $t=8\ell+16$ and $\delta=2^{-t}$, and compute
only the following rational quantities:
\begin{align*}
 S&=\operatorname{SQ}_t(X^2+Y^2) &&\text{on }[1,2],\\
 a'&=X/S,\\
 S_1&=\operatorname{SQ}_t((1+a')/2),&
 S_2&=\operatorname{SQ}_t((1+S_1)/2),\\
 S_3&=\operatorname{SQ}_t((1+S_2)/2),&
 T_3&=\operatorname{SQ}_t((1-S_2)/2),\\
 q_0&=S_3+i\sigma T_3.
\end{align*}
The last four square-root inputs lie in $[0,1]$. Because $S\ge1$,
$|a'-a|\le\delta$. Equation~\eqref{eq:sqrt-holder}, including the bisection
error at each stage, gives successive errors at most
$2\delta^{1/2}$, $2\delta^{1/4}$, and $2\delta^{1/8}$.
Both final coordinates have the last bound. Hence
\begin{equation}
 |q_0-v|\le4\delta^{1/8}=2^{-\ell},
 \qquad |q_0-1|\le5/8.
 \label{eq:q-close}
\end{equation}
The fractional powers in these error estimates are shorthand for repeated
square roots of the dyadic bound; the algorithm itself only performed the
five fixed-length rational bisections.

\subsection{A finite binomial identity}

Let $N\ge1$ and $\alpha_0=1/k$, and define rational numbers
\[
 c_0=1,\qquad c_{j+1}=c_j\frac{\alpha_0-j}{j+1},
 \qquad P_N(T)=\sum_{j=0}^N c_jT^j.
\]
Induction gives $|c_j|\le1$. A coefficient calculation gives
\begin{equation}
 (1+T)P_N'(T)-\alpha_0P_N(T)=(N-\alpha_0)c_NT^N.
 \label{eq:binomial-ode}
\end{equation}
Set $H(T)=P_N(T)^k=\sum_j d_jT^j$.
Since $k\alpha_0=1$, equation~\eqref{eq:binomial-ode} implies that
$(1+T)H'-H$ is divisible by $T^N$. Thus, for $j<N$,
\[
 (j+1)d_{j+1}=(1-j)d_j.
\]
It follows that $d_0=d_1=1$ and $d_j=0$ for $2\le j\le N$.
Therefore $H-(1+T)$ is divisible by $T^{N+1}$.

For $|T|\le5/8$, finite coefficient majorization now gives
\begin{align*}
 |P_N(T)^k-(1+T)|
 &\le (5/6)^{N+1}\sum_j|d_j|(3/4)^j\\
 &\le (5/6)^{N+1}
       \left(\sum_{j=0}^N|c_j|(3/4)^j\right)^k
 \le 4^k(5/6)^{N+1}.
\end{align*}
Take $N=8k+4\ell$. Since $(5/6)^4<1/2$, the last quantity is at most
$2^{-\ell}$. This argument uses finite polynomials and finite sums only;
it is not an invocation of an infinite binomial series.

\subsection{The output and its exact rational certificate}

Put $\eta=\varepsilon/(k+1)$. Choose
\[
 \ell=12+\bl\bigl(\ceil{(k+1)/\varepsilon}\bigr),
 \qquad N=8k+4\ell.
\]
For positive integers $m$, $\bl(m)$ is their binary digit count, so
$2^{\bl(m)}>m$. These choices ensure $\ell\ge3$ and
$2048\,2^{-\ell}\le\eta$.
Compute $q_0$ above and set
\begin{equation}
 u=P_N(q_0-1)^8,\qquad w=\frac{u}{1+\eta}.
 \label{eq:direction-output}
\end{equation}
Write $e=2^{-\ell}$ and $b=P_N(q_0-1)^k$.
Then $|b-q_0|\le e$, $|q_0-v|\le e$, and $|b|,|v|\le2$.
Equation~\eqref{eq:power-diff} for exponent $8$ gives
\[
 |u^k-c/|c||=|b^8-v^8|\le8\cdot2^7\cdot2e
                  =2048e\le\eta.
\]
Thus $|u|^k\le1+\eta$, and consequently $|u|\le1+\eta$.
For the latter implication, split the decidable comparison $|u|^2\le1$;
if it fails, $|u|\le|u|^k$. Hence $|w|\le1$.
Finally, with $r=(1+\eta)^{-1}$,
\[
 |w^k-c/|c||\le\eta r^k+(1-r^k)
       \le\eta+k(1-r)\le(k+1)\eta=\varepsilon.
\]
This proves Lemma~\ref{lem:direction}.

For clarity about arithmetization, put
\[
 A_c=\nn c,\quad b_c=\nn{w^k}+1-\varepsilon^2,
 \quad d_c=2\re(w^k\overline c).
\]
Here $b_c>0$. The output satisfies the \emph{purely rational} certificate
\begin{equation}
 \nn w\le1,\qquad d_c\ge0,\qquad d_c^2\ge A_cb_c^2.
 \label{eq:direction-certificate}
\end{equation}
Indeed the directional error bound is equivalent to
$b_c\sqrt{A_c}\le d_c$. Squaring positive sides gives
\eqref{eq:direction-certificate}; the finite bisection reasoning of
Section~\ref{sec:finite} removes the remaining radical notation.
The companion reference program checks precisely this certificate.

\section{Uniform descent for a rational polynomial}
\label{sec:descent}

\begin{lemma}[Rational descent with a radius certificate]
\label{lem:descent}
Let $n\ge1$, let $P\in\GQ[X]$ have degree at most $n$ and at least one
nonzero positive-degree coefficient, and let $z\in\GQ$. Put
\[
 D=2^n4^{2n^2},\qquad q=1-\frac1{4D}.
\]
A primitive recursive procedure returns $\Delta\in\GQ$ and rational
$r_0\ge0$ such that, writing $P(z+W)=\sum_{i=0}^n b_iW^i$,
\begin{align}
 |P(z+\Delta)|&\le q|P(z)|,\label{eq:descent-contraction}\\
 |\Delta|&\le r_0,\label{eq:descent-radius}\\
 |b_i|r_0^i&\le|b_0| \quad(1\le i\le n).
 \label{eq:radius-cert}
\end{align}
Every comparison in the procedure is rational and every search is bounded.
\end{lemma}

\begin{proof}
Translation uses only the finite binomial identity. If $b_0=0$, return
$(\Delta,r_0)=(0,0)$. Otherwise let $A=|b_0|>0$ and
\[
 M(r)=\max_{1\le i\le n}|b_i|r^i\quad(r>0).
\]
There is a nonzero $b_i$ with $i\ge1$, because inverse translation recovers
the positive-degree coefficients of $P$ from these coefficients.

\paragraph{A balanced starting radius.}
Let $b\ge1$ be the largest binary digit count of the absolute numerators and
positive denominators of all components of all $b_i$. Thus every coefficient
has norm at most $2^{b+1}$, and every nonzero coefficient has norm at least
$2^{-b}$. Set $E=2b+3$. At $r=2^{-E}$,
\[
 M(r)\le2^{b+1-E}=2^{-b-2}<A.
\]
At $r=2^E$, a nonzero positive-degree coefficient gives
$M(r)\ge2^{E-b}>A$. Choose the largest integer $e\in[-E,E]$ such that
$M(2^e)\le A$, and set $r_0=2^e$. This is a bounded search and $e<E$.
Since $M(2r_0)>A$ and $M(2r_0)\le2^nM(r_0)$,
\begin{equation}
 A/2^n<M(r_0)\le A.
 \label{eq:balanced}
\end{equation}
The right inequality proves \eqref{eq:radius-cert}. Tests use
$\nn{b_i}r^{2i}\le\nn{b_0}$; no square roots are computed.

\paragraph{A term dominant at three consecutive scales.}
For $j=0,\ldots,2n$, put $r_j=4^{-j}r_0$, and let $k_j$ be the least
index maximizing $|b_i|r_j^i$. The sequence $k_j$ is nonincreasing.
Indeed, if $k$ maximizes at $r$, an index $i>k$ has at $r/4$ a value at
most $4^{-i}|b_k|r^k<4^{-k}|b_k|r^k$, and cannot maximize there.
The strict inequality holds because the maximum is positive.

Some $j\in\{1,\ldots,2n-1\}$ satisfies
$k_{j-1}=k_j=k_{j+1}$. To see this by finite arithmetic, examine the odd
centers. If all $n$ odd-centered triples failed, monotonicity would give
$k_{2s+2}<k_{2s}$ for $s=0,\ldots,n-1$, hence
$k_{2n}\le k_0-n\le0$, a contradiction. Bounded search finds such a $j$.

Set $r=r_j$, $k=k_j$, and $V=|b_k|r^k$. Dominance at $r/4$ gives, for
$i<k$, $|b_i|r^i\le4^{-(k-i)}V$; dominance at $4r$ gives the analogous
$4^{-(i-k)}V$ bound for $i>k$. Therefore
\begin{equation}
 \sum_{\substack{1\le i\le n\\i\ne k}}|b_i|r^i\le\frac23V.
 \label{eq:other-terms}
\end{equation}
Each of the two finite geometric sums is at most $1/3$.
Also $0<V\le A$. A term maximizing at $r_0$ loses a factor of at most
$4^{jn}$ by radius $r_j$. Since $j<2n$, \eqref{eq:balanced} gives
\begin{equation}
 V>A/(2^n4^{2n^2})=A/D.
 \label{eq:dominant-lower}
\end{equation}

\paragraph{The direction and the decrease.}
Apply Lemma~\ref{lem:direction} with
$c=-b_0\overline{b_k}\ne0$, exponent $k$, and $\varepsilon=1/16$.
Its unit target is
$u=-b_0\overline{b_k}/(A|b_k|)$.
Return $\Delta=rw$, where $w$ is the rational direction output.
Then $|\Delta|\le r\le r_0$ and
$b_kr^ku=-b_0V/A$. Using the weaker bound $1/12$ for the direction error,
\[
 |b_0+b_k\Delta^k|\le A-V+V/12.
\]
The other terms have total norm at most $2V/3$ by
\eqref{eq:other-terms} and $|w|\le1$. Hence
\[
 |P(z+\Delta)|\le A-\frac14V
       \le\left(1-\frac1{4D}\right)A.
\]
This proves all claims. The actual $1/16$ direction estimate leaves a
strict margin in this calculation.
\end{proof}

The lemma does not rely on monicity or on the existence of any root. Its
conclusion is a finite rational statement. It is the key replacement for
an appeal to a minimum of $|P|$ on a compact set.

\section{A bound on travel from any nonconstant coefficient}
\label{sec:travel}

The residual contraction alone does not make the iterates a Cauchy sequence.
We next extract a step bound from the supplied nonconstant coefficient. This
also avoids any need to determine the exact degree.

\begin{lemma}[Step bound]
\label{lem:step-bound}
Suppose $P(X)=\sum_{i=0}^n p_iX^i$, $|p_k|\ge\lambda>0$ for
$1\le k\le n$, and $|z|\le U-1$ for rational $U\ge1$.
Let the descent lemma return $(\Delta,r_0)$. If
\[
 |P(z)|\le BQ^j,\qquad B\ge1,\qquad 0<\rho<1,
 \qquad Q\le\rho^n,
\]
and $C=Bn2^n/\lambda\ge1$, then
\[
 |\Delta|\le r_0\le UC\rho^j.
\]
\end{lemma}

\begin{proof}
The case $P(z)=0$ has zero step. Otherwise use inverse translation:
\begin{equation}
 p_k=\sum_{i=k}^n\binom ik b_i(-z)^{i-k}.
 \label{eq:inverse-translation}
\end{equation}
There are at most $n$ terms, $\binom ik\le2^n$, and
$|z|^{i-k}\le U^i$. Put $F=n2^n$. The triangle inequality gives
\begin{equation}
 |p_k|\le F\max_{1\le i\le n}|b_i|U^i.
 \label{eq:translate-lower}
\end{equation}
Equivalently, a completely rational consequence of finite Cauchy--Schwarz is
\begin{equation}
 \nn{p_k}\le F^2\max_{1\le i\le n}\nn{b_i}U^{2i}.
 \label{eq:translate-lower-squared}
\end{equation}
Choose an index maximizing the displayed rational expression. It has
$|b_i|U^i\ge\lambda/F>0$. The radius certificate gives
\[
 (r_0/U)^i\le(F/\lambda)|P(z)|
       \le C\rho^{nj}\le(C\rho^j)^i.
\]
The last inequality uses $i\le n$, $C\ge1$, and $0<\rho<1$.
Positive integer powers reflect order on nonnegative rationals, proving
the claim. Alternatively, use \eqref{eq:translate-lower-squared} and the
squares of the radius and residual bounds to obtain the same rational
comparison. No operation taking an $i$th root is needed here.
\end{proof}

\section{All constants and the primitive recursive iteration}
\label{sec:iteration}

Return to the input of Theorem~\ref{thm:main}. Define the positive integers
and positive rationals
\begin{align}
 B&=2+\ceil{|\re a_0(0)|+|\im a_0(0)|},
 &D&=2^n4^{2n^2},\nonumber\\
 q&=1-\frac1{4D},&Q&=1-\frac1{8D},\nonumber\\
 K&=8nD,&\rho&=1-\frac1K,\nonumber\\
 C&=Bn2^{n+h+1}.&&
 \label{eq:constants}
\end{align}
In particular $B\ge2$, $C\ge B$, and $K\ge2$.
Finite Bernoulli arithmetic gives
\begin{equation}
 \rho^n\ge1-n/K=Q.
 \label{eq:rho-q}
\end{equation}

The following choices make every bound explicit; binary digit counts are
primitive recursive:
\begin{align}
 A_0&=\bl(2CK),&T&=KA_0,&R&=2(1+C)^T,
 \label{eq:radius-global}\\
 L&=\max\{h+2,\bl(16D(n+1)R^n)\},&&&
 \label{eq:precision-global}\\
 H&=\bl(RCK).&&&
 \label{eq:H}
\end{align}
Thus
\[
 2^{A_0}\ge2CK,\quad 2^L\ge16D(n+1)R^n,
 \quad L\ge h+2,\quad 2^H\ge RCK.
\]
The numbers are intentionally generous. Small asymptotic complexity is not
needed for primitive recursiveness.

Set
\begin{equation}
 N_j=j+L,\qquad P_j(X)=\sum_{i=0}^n a_i(N_j)X^i.
 \label{eq:sample-polynomial}
\end{equation}
Every $P_j$ is a rational nonconstant polynomial, since
$|a_k(N_j)|\ge\lambda=2^{-h-1}$ by \eqref{eq:apart-tail}.
Starting at $z_0=0$, define
\begin{equation}
 (\Delta_j,r_j^*)=\operatorname{Descent}(P_j,z_j),
 \qquad z_{j+1}=z_j+\Delta_j.
 \label{eq:iteration}
\end{equation}
The superscript on $r_j^*$ distinguishes the returned starting radius from
the local scale indices used in the proof of the descent lemma.
Equation~\eqref{eq:iteration} is one ordinary primitive recursion on $j$,
with a primitive recursive body. The coefficient program is fixed, so calls
to it are ordinary primitive recursive compositions.

For the proof, also define the rational sequence
\begin{equation}
 U_0=1,\qquad U_{j+1}=U_j(1+C\rho^j).
 \label{eq:U}
\end{equation}
It is an auxiliary bound, not a new infinite object assumed to exist.

\subsection{The auxiliary products are bounded}

First observe the useful block inequality
\begin{equation}
 \rho^K\le\frac12.
 \label{eq:block}
\end{equation}
Indeed,
\[
 \rho^{-K}=\left(1+\frac1{K-1}\right)^K
     \ge1+\frac K{K-1}\ge2,
\]
by finite Bernoulli arithmetic; all denominators are positive.
For every finite $\ell\ge T$, the geometric sum gives
\begin{equation}
 \sum_{j=T}^{\ell-1}C\rho^j
       \le CK\rho^T\le CK2^{-A_0}\le\frac12.
 \label{eq:tail-sum}
\end{equation}

For nonnegative rational $x_1,\ldots,x_m$ with $S=\sum x_i<1$,
\begin{equation}
 \prod_{i=1}^m(1+x_i)\le\frac1{1-S}.
 \label{eq:finite-product}
\end{equation}
To prove this, induct on the finite list. If the previous sum is $s$ and
the next term is $x$, then
\[
 (1+x)(1-s-x)=1-s-sx-x^2\le1-s.
\]
Multiplying the previous product bound by $1+x$ proves the new one.
Applying \eqref{eq:finite-product} to \eqref{eq:tail-sum} bounds every
tail product by $2$. The first $T$ factors are each at most $1+C$.
Consequently, for every $j$, including $j<T$,
\begin{equation}
 1\le U_j\le2(1+C)^T=R.
 \label{eq:U-bound}
\end{equation}
This is a theorem about finite products only. No existence or convergence
of an infinite product has been used.

\section{The rational induction invariant}
\label{sec:invariant}

\begin{proposition}
\label{prop:invariant}
For every $j\in\N$,
\begin{align}
 |P_j(z_j)|&\le E_j:=BQ^j,\label{eq:residual-invariant}\\
 |z_j|&\le U_j-1<R,\label{eq:position-invariant}\\
 |\Delta_j|&\le U_jC\rho^j\le RC\rho^j.
 \label{eq:step-invariant}
\end{align}
\end{proposition}

\begin{proof}
We induct on the first two inequalities. The third is a consequence used
in the induction step. At $j=0$, $z_0=0=U_0-1$, and the name hypothesis
gives
\[
 |P_0(0)|=|a_0(L)|\le|a_0(0)|+1+2^{-L}\le B.
\]

Suppose the first two inequalities hold at $j$. The sampled coefficient
$a_k(N_j)$ has norm at least $\lambda=2^{-h-1}$.
Equations~\eqref{eq:rho-q} and \eqref{eq:constants} make
Lemma~\ref{lem:step-bound} applicable with $U=U_j$ and
$C=Bn2^n/\lambda$. It gives \eqref{eq:step-invariant}. Therefore
\[
 |z_{j+1}|\le U_j-1+U_jC\rho^j=U_{j+1}-1<R.
\]

It remains to control the change of coefficient approximations.
For any rational $x$ with $|x|\le R$, since $R\ge1$ and $N_{j+1}=N_j+1$,
\begin{align}
 |P_{j+1}(x)-P_j(x)|
 &\le\frac32(n+1)R^n2^{-N_j}\nonumber\\
 &\le\frac3{32D}2^{-j}
 \le\frac1{8D}BQ^j.
 \label{eq:drift}
\end{align}
The first inequality uses the coefficient-name hypothesis, the second
uses the choice of $L$, and the last uses $B\ge1$, $Q\ge1/2$, and
$3/32\le1/8$. The descent contraction at $z_j$, followed by
\eqref{eq:drift} at the already bounded point $z_{j+1}$, yields
\[
 |P_{j+1}(z_{j+1})|
 \le qE_j+\frac1{8D}E_j
 =QE_j=E_{j+1}.
\]
This completes the induction.
\end{proof}

The important logical feature is the use of $P_j(z_j)$, which is an actual
Gaussian rational. The invariant is \emph{not} an assertion about the value
of a completed complex polynomial at every possible precision. Its squared
form is simply
\begin{equation}
 \nn{P_j(z_j)}\le B^2Q^{2j},\qquad
 \nn{z_j}\le(U_j-1)^2.
 \label{eq:qf-invariant}
\end{equation}
These are quantifier-free numerical comparisons once the primitive
recursive operations have been named. The global name assumptions remain
in the context; they are instantiated at finite indices during the proof.
Section~\ref{sec:ipra-audit} gives an entirely finite-table formulation.

\section{The exact root and its explicit modulus}
\label{sec:root}

For $\ell\ge j$, finite telescoping and \eqref{eq:step-invariant} give
\begin{equation}
 |z_\ell-z_j|\le RC\sum_{r=j}^{\ell-1}\rho^r
                         \le RCK\rho^j.
 \label{eq:iterate-tail}
\end{equation}
Define the promised functions
\begin{equation}
 J(t)=K(t+H+2),\qquad
 \alpha(t)=z_{J(t)},\qquad M(t)=J(t)+L.
 \label{eq:root-output}
\end{equation}
All constants are fixed by the input and all three functions are primitive
recursive. For $u\ge t$, \eqref{eq:block}, \eqref{eq:H}, and
\eqref{eq:iterate-tail} give
\begin{equation}
 |\alpha(u)-\alpha(t)|\le RCK\rho^{J(t)}
      \le RCK2^{-t-H-2}\le2^{-t-2}.
 \label{eq:root-cauchy}
\end{equation}
This proves the fast Cauchy assertion, with a displayed modulus.

Let $\beta(t)=P_{J(t)}(\alpha(t))$. Since $C\ge B$, $R,K\ge1$, and
$Q\le\rho^n\le\rho$, the residual invariant gives
\begin{equation}
 |\beta(t)|\le BQ^{J(t)}
     \le B\rho^{J(t)}\le2^{-t-2}.
 \label{eq:beta-zero}
\end{equation}
For every $s\ge M(t)=J(t)+L$, the coefficient-name hypothesis gives
\begin{align*}
 \left|\sum_{i=0}^n a_i(s)\alpha(t)^i-\beta(t)\right|
 &\le2(n+1)R^n2^{-J(t)-L}\\
 &\le\frac1{8D}2^{-J(t)}\le2^{-t-3}.
\end{align*}
Adding \eqref{eq:beta-zero} proves \eqref{eq:main-root}, in fact with
$3\cdot2^{-t-3}$ in place of $2^{-t}$. This establishes both assertions of
Theorem~\ref{thm:main} using only finite rational estimates.

\subsection{Why this is an exact zero, not just approximate solvability}

A complex number in this representation is a fast Cauchy name up to the
usual vanishing-error equivalence. The object $\alpha$ is the specific
primitive recursive name \eqref{eq:root-output}. We do not infer its
existence by extracting a subsequence or choosing compatible approximate
roots.

The construction acts on the coefficient names themselves. Equivalent
names can cause it to select different valid roots. Constructive root
existence does not require an extensional single-valued selection of a
root from every complex polynomial, and no such selection is asserted.

To connect the residual calculation explicitly with ordinary name
arithmetic, set
\[
 V(t)=\sum_{i=0}^n a_i(t)\alpha(t)^i,
 \qquad S=(n+1)R^n.
\]
The input name estimate and $M(t)\ge t$ imply
\[
 |V(t)-\beta(t)|\le2S2^{-t},\qquad
 |V(t)|\le(2S+1/4)2^{-t}.
\]
Choose a fixed integer $d$ with $2^d\ge4S+1$.
Then $|V(t+d)|\le2^{-t-1}$. Thus this normalized evaluation sequence
represents zero and satisfies the usual fast Cauchy condition as well.

For extensionality, a bound for all coefficient samples is
\[
 B_* =2+\ceil{\max_{0\le i\le n}
        (|\re a_i(0)|+|\im a_i(0)|)}.
\]
For $|x|,|y|\le R$, finite identity~\eqref{eq:power-diff} gives
\begin{equation}
 \left|\sum_{i=0}^n a_i(s)(x^i-y^i)\right|
       \le B_*\left(\sum_{i=1}^n iR^{i-1}\right)|x-y|.
 \label{eq:poly-lipschitz}
\end{equation}
Changing coefficient samples changes a value by at most the corresponding
coefficient error times $\sum_{i=0}^nR^i\le S$.
These two explicit estimates show that replacing either the coefficient
names or the root name by equivalent names preserves the evaluation class.
They also give a fixed primitive recursive precision shift for any of the
standard conventions for addition and multiplication of fast names.
Consequently \eqref{eq:beta-zero} is exactly $p(\alpha)=0$ in the
represented complex numbers. No real or complex completeness principle is
being inserted at this final step.

\section{Why the proof belongs to \texorpdfstring{$\ipra$}{iPRA}}
\label{sec:ipra-audit}

The size of the numerical bounds is immaterial to membership in the
primitive recursive class. Their \emph{form} and the induction formulas
are what matter. We make both explicit.

\subsection{Total primitive recursive definitions}

Every object used by the root program is obtained by the following
operations: natural and integer arithmetic; rational arithmetic with
positive denominators; finite tuple coding; bounded maxima and bounded
search; binary digit count; finite exponentiation; finite polynomial
operations; fixed-length rational bisection; and one primitive recursion
forming the iterates. If reduced fractions are used, gcd can itself be
defined by a bounded search. Reduced fractions are not required for the
arithmetical proof.

The local radius search has at most $2E+1$ candidate exponents. The scale
search uses $2n+1$ maxima and at most $2n-1$ possible centers. The direction
routine has five bisection loops of length $8\ell+16$, a binomial recurrence
of length $8k+4\ell$, and finite Horner evaluation. The complete root program
runs the local recursion exactly $J(t)$ times. Dependence of a local loop
bound on its current rational input is harmless: it is a primitive
recursive function forming part of the fixed recursion body.

For total definitions on all natural codes, assign a fixed default output
to malformed rational encodings, $n=0$, a missing nonconstant coefficient,
or a failed bounded search. The hypotheses and lemmas prove that these
default branches are never used on valid inputs. This supplies total PR
function symbols without presupposing their correctness.

\subsection{The finite-table theorem}

Fix a finite iteration horizon $m$. Form a finite table containing
$a_i(0)$, $a_i(h+2)$, and $a_i(L+j)$ for $0\le i\le n$ and $0\le j\le m$.
The table conditions needed for Proposition~\ref{prop:invariant} are only:
\begin{enumerate}
\item the finite apartness comparison at $(k,h+2)$;
\item the comparisons of that row with the rows $L+j$;
\item the comparison of the constant coefficient at $0$ with the row $L$;
\item the comparisons of adjacent rows $L+j$ and $L+j+1$.
\end{enumerate}
All are rational inequalities. Their conjunction over finite index ranges
is a primitive recursive Boolean $\mathsf{GoodTable}$.
Keep this table and $m$ as parameters. Induct on a numerical assertion of
the form
\[
 j\le m\ \longrightarrow\
 (\mathsf{GoodTable}=1\ \longrightarrow\ \mathsf{Inv}(j)=1),
\]
where $\mathsf{Inv}$ checks \eqref{eq:qf-invariant} and any accumulated
finite bounds. This is a quantifier-free formula in the full primitive
recursive language. The descent lemma and the preceding arithmetic lemmas
supply the step. The finite tail estimates and the arbitrary-precision
residual comparison are treated with extra finite indices as parameters.

When the table is generated by the fixed PR function $A$, the global
name hypothesis proves these finite table conditions by instantiation and
bounded Boolean-fold induction. The hypotheses themselves need not be
quantifier-free: they remain assumptions in the context. It is the
\emph{formula to which induction is applied} that must be quantifier-free.
After the finite proof, ordinary intuitionistic universal introduction
quantifies over the arbitrary horizons and precisions.

\subsection{An induction and logical-principle audit}

\begin{longtable}{@{}>{\raggedright\arraybackslash}p{0.24\textwidth}>{\raggedright\arraybackslash}p{0.39\textwidth}>{\raggedright\arraybackslash}p{0.29\textwidth}@{}}
\toprule
Part & Finite invariant or calculation & Induction parameter \\
\midrule
\endhead
Square-root bounds & Nonnegative rational endpoints, squared bracketing,
interval length & Bisection count \\
Binomial truncation & Coefficient recurrence and vanishing coefficients
through degree $N$ & Finite coefficient index \\
Radius and scale selection & Rational squared maxima; decreasing
maximizer indices & Bounded scan index \\
Direction output & Certificate \eqref{eq:direction-certificate} after
finite arithmetic and bisection bounds & Fixed finite precision loops \\
Translation bound & Finite polynomial identity and squared
Cauchy--Schwarz estimate & Degree or list index \\
Auxiliary product & Rational partial sum and partial product inequalities
& Product length \\
Iteration & Two rational inequalities \eqref{eq:qf-invariant} under a
fixed good-table condition & Iteration index \\
Root name & Finite telescoping and finite geometric bounds
& Finite interval length \\
\bottomrule
\end{longtable}

The proof uses no excluded middle for real equality, unbounded minimization,
compactness, infinite path principle, general intermediate value theorem,
countable choice, Markov principle, or induction on an unbounded
existential or universal formula. Rational case distinctions are justified
by decidability derived in arithmetic. The few contradiction arguments
prove decidable rational or finite combinatorial assertions and are
intuitionistically valid.

It would be insufficient to say only that a numerical algorithm appears
computable, or that Lean accepts a theorem about standard $\C$.
The finite bounds and the quantifier-free invariants above are the
additional content establishing the $\ipra$ claim.

\section{Familiar corollaries and the exact scope}
\label{sec:corollaries}

\subsection{Monic and specified-degree polynomials}

For a monic polynomial, the top coefficient is the constant name $1$, so
$(k,h)=(n,0)$ supplies the required certificate. All lower coefficients may
be arbitrary fixed PR fast names. A polynomial with a supplied apart
leading coefficient is already covered by the general theorem, without
normalization.

There is a shorter bound in the monic case. The local radius certificate
then gives $r_0^n\le|P_j(z_j)|$, so the step bound improves to
$|\Delta_j|\le B\rho^j$. One may take $R=BK+1$ and omit the auxiliary
product argument. The general construction was retained in the main proof
to handle any positively nonconstant coefficient.

\subsection{Complete factorization for every fixed degree}

Suppose a monic polynomial has fixed degree $n$ and let $\alpha$ be the
constructed root. Synthetic division over the represented complex numbers
is finite algebra. Define
\[
 q_{n-1}=1,\qquad q_{j-1}=a_j+\alpha q_j
       \quad(j=n-1,n-2,\ldots,1).
\]
Then direct coefficient comparison gives
\[
 p(X)=(X-\alpha)\sum_{j=0}^{n-1}q_jX^j+p(\alpha)
      =(X-\alpha)\sum_{j=0}^{n-1}q_jX^j.
\]
The quotient is monic. The finite arithmetic of PR names, with the explicit
bounded-domain precision shifts from \eqref{eq:poly-lipschitz}, makes each
$q_j$ a PR name. Repeat the root construction $n-1$ times. For each fixed
$n$ this is a finite composition of the preceding program constructions
and yields
\[
 p(X)=\prod_{j=1}^n(X-\alpha_j).
\]
With a specified apart leading coefficient, division by that coefficient
first gives the usual leading scalar times the product. Inversion has a
PR modulus: if $|a|\ge2^{-h}$, sufficiently accurate samples have norm
at least $2^{-h-1}$, and
\[
 |a(s)^{-1}-a^{-1}|\le2^{2h+1}|a(s)-a|.
\]
A fixed precision shift therefore suffices.

This paragraph claims a \emph{fixed-degree factorization schema}. It does
not assert, without a further compiler-size argument, that evaluating the
nested factorization uniformly over arbitrary degree and arbitrary PR
program codes is itself PR. The single-root theorem proved above is
already uniform in the numerical degree bound and is the FTA formulation
requested here.

\subsection{What has and has not been proved}

The mathematical result is an explicit arithmetized FTA for PR Cauchy
coefficients, including the full positive-nonconstancy formulation. It is
not a claim that first-order arithmetic literally quantifies over all
uncountably many ordinary complex numbers. It is not an assertion that weak
negated equality supplies positive apartness. It is also not yet a
machine-checked Lean derivation. These distinctions specify the theorem's
content; none is filled by an extra axiom.

\section{A concrete path to complete Lean formalization}
\label{sec:lean}

\subsection{Inspected repository foundation}

The repository inspection was pinned to ProveIt commit
\begin{center}
\texttt{b03925cccc76d55b9dc2a00864a5e0544b64aa7f}.
\end{center}
The source uses Lean and mathlib \texttt{v4.32.0}. These were inspected, not
rebuilt in this environment. The small \texttt{ConciseIPRA} foundation is
the recommended dependency. In namespace \texttt{LeanProofs.ConciseIPRA},
the relevant actual names are:

\begin{center}
\begin{tabular}{@{}ll@{}}\toprule
Repository object & Role \\
\midrule
\texttt{PR k} & Primitive recursive function syntax \\
\texttt{PTm n} & Arithmetic terms with $n$ variables \\
\texttt{PRAAx} & Defining equations and QF induction \\
\texttt{Prf Ax Gamma phi} & Intuitionistic object-theory derivability \\
\texttt{PRA.qfInd} & Derived quantifier-free induction rule \\
\texttt{PRA.precZero}, \texttt{PRA.precSucc} & Recursion equations \\
\texttt{pra\_iff\_teq}, \texttt{pra\_iff\_native} & Transfers to concise presentations \\
\bottomrule\end{tabular}
\end{center}

The decisive endpoint, after closing the schema instance, is
\begin{lstlisting}
Prf PRAAx [] sentence
\end{lstlisting}
where \texttt{sentence} is the actual arithmetical encoding of the input
hypotheses implying the two output bounds. An ordinary Lean proof of
\texttt{Complex} algebraic closedness is a different result.

The inspected DiophantineConsistency development supplies another
embedded iPRA presentation and finite PR syntax, but is unnecessary for
this proof. The inspected Busy Beaver/ZFC entry points concern machine
and consistency constructions; no theorem from them is used as an FTA
premise. Avoiding those larger dependencies makes the logical strength of
the present proof easier to audit.

\subsection{Modules, exact obligations, and proof strategies}

The included \texttt{LEAN\_PLAN.md} gives the implementation plan in more
detail. The essential milestones are as follows.

\begin{enumerate}
\item \textbf{Numerical codes.} Define total integer, rational, Gaussian
rational, and finite-list operations as \texttt{PR} codes. Define the
Boolean predicates for well-formedness, rational order, and squared-norm
bounds. Prove their equations and basic arithmetic internally using
\texttt{PRA.precZero}, \texttt{PRA.precSucc}, and \texttt{PRA.qfInd}.

\item \textbf{Finite algebra.} Prove polynomial translation and inverse
translation, \eqref{eq:power-diff}, finite sums and products, binomial
identities, and squared Cauchy--Schwarz. Every sequence is explicitly coded,
and every induction invariant is a Boolean-coded finite check.

\item \textbf{Direction.} Implement the five bisections, rational binomial
recurrence, and Horner evaluation. Prove the polynomial cancellation in
\eqref{eq:binomial-ode} coefficient by coefficient, then the finite
weighted coefficient bound. End with the rational certificate
\eqref{eq:direction-certificate}; do not import FTA to choose a root.

\item \textbf{Descent.} Define the bounded radius search and scale maxima.
Prove their certificates, monotonicity of the maximizing indices, the
odd-center triple argument, and \eqref{eq:descent-contraction}--
\eqref{eq:radius-cert} as squared-rational formulas. Include the
$b_0=0$ branch explicitly.

\item \textbf{Iteration.} Define all constants as PR programs. Prove
\eqref{eq:translate-lower-squared}, the product bound, and the finite-table
version of \eqref{eq:qf-invariant}. This is the central internal
\texttt{PRA.qfInd} application.

\item \textbf{Name theorem.} Prove finite telescoping, substitute
$J(t)=K(t+H+2)$, and derive the two rational output assertions. Generate
the finite coefficient table from a fixed \texttt{PR} coefficient program.
Use first-order quantifier introduction to obtain each closed schema
instance. Produce a term of \texttt{Prf PRAAx [] sentence}.

\item \textbf{Optional presentation transfer.} Only after obtaining that
object-level derivation, use \texttt{pra\_iff\_teq} or
\texttt{pra\_iff\_native} to transfer it to the concise presentation.
An external interpretation by fast Cauchy names then gives the familiar
complex-root reading.
\end{enumerate}

Lean's ordinary rational arithmetic can be used first as a development and
testing layer. It does not replace the internal proof. Reification or a
small checked proof-producing arithmetic library can share much of the
work: each numerical transformation should return an object-theory proof
of its corresponding equation or implication. Neither \texttt{norm\_num}
on a single example nor a general theorem in Lean's metatheory alone
establishes the universal iPRA derivation.

\subsection{Delivered formalization status}

\texttt{lean/Specification.lean} is explicitly labeled an uncompiled
interface and blueprint. It contains specifications and types for the
eventual certificates, with no claimed completed FTA proof. There was no
Lean executable in the working environment, so no claim is made that the
file has passed Lean's elaborator. The archive does not turn missing
formal proofs into axioms asserting the theorem.

The handwritten proof above supplies the mathematical content of all the
listed obligations, including the direction construction, explicit
precision shifts, and finite induction audit. The remaining work is
formal implementation of the numerical coding library and its object-level
derivations, followed by their assembly. This is a substantial formalization
project, but it does not leave a missing mathematical existence lemma or
an unspecified termination argument.

\section{Reproducible artifacts and validation}

The archive contains this LaTeX source and its compiled PDF; the Lean
specification and module plan; exact rational implementations of the
direction routine and one descent step; and recorded test results.
From the \texttt{code} directory, run
\begin{lstlisting}
python3 test_direction.py
python3 test_descent.py
\end{lstlisting}
The construction and checks use Python's exact rational arithmetic, not
floating-point approximations. Direction tests cover several quadrants,
exponents $1,2,3,5$, and a target with irrational absolute value; their
certificate is the rational inequality
\eqref{eq:direction-certificate}. Descent tests include a quadratic, a
zero highest coefficient, an already reached root, and translation away
from zero. They check the squared residual contraction and the returned
radius certificates exactly.

These computations check the reference implementations on the listed
inputs. They are not a proof of the universal theorem and are not a Lean
verification. The deliberately large worst-case schedule $J(t)$ is given
mathematically; the code does not claim to run that entire schedule as a
practical root solver. Its purpose is to make the finite constructions
concrete and independently inspectable.

\section*{Provenance}
\addcontentsline{toc}{section}{Provenance}

The local descent idea belongs to the Kneser constructive tradition.
Geuvers, Wiedijk, and Zwanenburg give a primary account of a constructive
FTA proof and its Coq formalization~\cite{kneser}. The bounded rational
direction construction, explicit constants, general coefficient travel
bound, and finite-table iPRA organization presented here are developed
in full above. The proof is not obtained by importing the Coq theorem or
by appealing to ordinary Lean/mathlib FTA.

The user's cited Wikipedia page~\cite{wikipedia} fixes the usual theorem
being addressed. The ProveIt sources~\cite{proveit-pr,proveit-deduction,
proveit-bi,proveit-native} fix the recommended object-theory and translation
interfaces. No consistency assertion from the Busy Beaver or
DiophantineConsistency developments is a premise of the argument.

\begin{thebibliography}{9}
\raggedright
\bibitem{kneser}
Herman Geuvers, Freek Wiedijk, and Jan Zwanenburg,
\emph{A Constructive Proof of the Fundamental Theorem of Algebra without
using the Rationals}, author-hosted paper.
\url{https://www.cs.ru.nl/~freek/pubs/kneser.pdf}.

\bibitem{wikipedia}
\emph{Fundamental theorem of algebra}, Wikipedia, the page specified in
the request.
\url{https://en.wikipedia.org/wiki/Fundamental_theorem_of_algebra}.

\bibitem{proveit-pr}
Vladimir Reshetnikov, ProveIt,
\path{Logic/PrimitiveRecursiveArithmetic/ConciseIPRA/Lean/ConciseIPRA/PR.lean},
commit \texttt{b03925cccc76d55b9dc2a00864a5e0544b64aa7f}.
\href{https://github.com/VladimirReshetnikov/ProveIt/blob/b03925cccc76d55b9dc2a00864a5e0544b64aa7f/Logic/PrimitiveRecursiveArithmetic/ConciseIPRA/Lean/ConciseIPRA/PR.lean}{Pinned source}.

\bibitem{proveit-deduction}
Same repository and commit, \texttt{ConciseIPRA/Deduction.lean}.
\href{https://github.com/VladimirReshetnikov/ProveIt/blob/b03925cccc76d55b9dc2a00864a5e0544b64aa7f/Logic/PrimitiveRecursiveArithmetic/ConciseIPRA/Lean/ConciseIPRA/Deduction.lean}{Pinned source}.

\bibitem{proveit-bi}
Same repository and commit, \texttt{ConciseIPRA/BiInterpretation.lean}.
\href{https://github.com/VladimirReshetnikov/ProveIt/blob/b03925cccc76d55b9dc2a00864a5e0544b64aa7f/Logic/PrimitiveRecursiveArithmetic/ConciseIPRA/Lean/ConciseIPRA/BiInterpretation.lean}{Pinned source}.

\bibitem{proveit-native}
Same repository and commit, \texttt{ConciseIPRA/NativeBiInterpretation.lean}.
\href{https://github.com/VladimirReshetnikov/ProveIt/blob/b03925cccc76d55b9dc2a00864a5e0544b64aa7f/Logic/PrimitiveRecursiveArithmetic/ConciseIPRA/Lean/ConciseIPRA/NativeBiInterpretation.lean}{Pinned source}.
\end{thebibliography}
\end{document}

